Fuzzing is nowadays a widely used approach for software testing and vulnerability assessment \cite{Zhu2022,Mallissery2023}.
    Its main strength derives from the fact that it can rapidly produce various inputs to be fed into a target executable.
    In general, a fuzzer can be divided into the \textbf{Input Generator}, which is responsible for providing inputs to the \textbf{Executor} that executes them.
    During the execution, the \textbf{Defect Monitor} checks whether the executed input causes a crash of the target.\\
    In particular, the Input Generator must choose which input to provide to the executor among all the inputs in the \textit{Input Space}.
    A subset of the \textit{Input Space} is represented by the \textit{Valid Input Space}.
    This is a subset of the input space that specifically contains all the inputs that are well formed with respect to the
    application being tested.
    For example, if the fuzzed target is a PDF reader, a valid input is a PDF file.
    A good Input Generator must choose inputs that are also valid inputs, because only by doing this is there a probability of finding a vulnerability \cite{Zhu2022}.
    Hence the need for a precise input selection and mutation strategy arises, in order to discover as many edge cases as possible.\\
    In order to do so, the Input Generator must obtain some information from the Executor.
    Depending on the amount of information the Input Generator has, fuzzing can be classified into
    \textit{White-box}, \textit{Black-box} or \textit{Grey-box}.
    In Black-box fuzzing, the fuzzer does not have any knowledge of the execution \cite{Miller1990}.
    In White-box fuzzing, instead, the source code must always be available, and techniques such as Symbolic and Concolic execution are used \cite{Godefroid2012}.
    In Grey-box fuzzing, instead, the inputs are chosen using, for example, \textbf{Code Coverage} information (i.e. the number of lines of code visited during execution of the target).
    This approach is called \textit{Coverage-Guided Grey-box fuzzing} and is based on prioritizing the inputs that discover new sections of code.
    Since the source code is not always available, Code Coverage is not always a viable Coverage Measure.
    In general, the coverage measures can be of different granularity \cite{CollAFL}.
    The three main options are:
    \begin{itemize}
        \item \textbf{Block Coverage}: Keeps track of the hit count for each \gls{BB}. This approach is used by some
            popular fuzzers such as libFuzzer or honggfuzz \cite{CollAFL}. This method is often too imprecise because
            it fails to capture the transitions between BBs.
        \item \textbf{Edge Coverage}: Keeps track of the hit count of every edge in the \gls{CFG} (i.e. the transition between two \gls{BB}s).
            This is the approach used by AFL and AFL++ \cite{aflplusplus}. It maintains overheads similar to Block Coverage but can
            track the transitions between \gls{BB}s and not only their execution, adding some context information.
            However, it lacks precision in keeping track of the order in which the edges are taken.
        \item \textbf{Path Coverage}: This is the most complete metric and keeps track of the complete path taken during the execution
            of a specific input, thus also keeping track of the order in which the edges are taken.
            The main drawback of this metric is that a single execution path could become very large, resulting in
            large runtime and memory overheads.
    \end{itemize}

    The best approach to Code Coverage, in terms of overhead and effectiveness, is \textbf{Edge Coverage} \cite{CollAFL}, which keeps track of which edges
    in the \gls{CFG} are taken during execution without the overheads incurred by keeping track of the entire execution path.
    Coverage information is obtained through the instrumentation of the target application.
    Depending on the availability of the source code, the instrumentation can be added in different ways:
    \begin{itemize}
        \item The source code is available: In this case the instrumentation is added at \textit{Compile Time}.
        \item The source code is not available: In this case there are two different techniques for instrumenting the application:
        \begin{itemize}
            \item \textbf{\gls{SBI}}: The instrumentation is added before the target is executed, without access to the
                source code, relying on Static Binary Rewriting techniques.
            \item \textbf{\gls{DBI}}: The instrumentation is added during program execution, relying on Dynamic Binary Rewriting techniques.
        \end{itemize}
    \end{itemize}
    Regardless of the availability of the source code, the \textit{Coverage-Guided Grey-box fuzzing} technique can be summarized
    as shown in \cref{fig:coverageguidedfuzzing}.
    \begin{figure}[ht]
        \caption{Generic diagram describing the Coverage-Guided Grey-box fuzzing workflow}
        \centering
        \includegraphics[width=\textwidth]{CoverageGuidedGrayBoxFuzzing.png}
        \label{fig:coverageguidedfuzzing}
    \end{figure}
    \\Specifically for \textbf{Edge Coverage}, \cref{fig:coverageguidedfuzzing} can be described as follows:
    \begin{enumerate}
        \item \textbf{Seed Corpus}: Represents the set of inputs on which mutations are computed.
            At the beginning it is composed of user-defined seeds, but during fuzzing this set can grow if a specific mutation discovers new edges not yet found.
        \item \textbf{Input Generator}: Composed of two main parts:
            \begin{itemize}
                \item \small \textbf{Seed Scheduling}: responsible for deciding which seed in the corpus to mutate.
                \item \small \textbf{Mutation Scheduling}: responsible for deciding which mutation or mutations to apply to the selected seed.
            \end{itemize}
        \item \textbf{Executor}: The component that actually executes the target binary.
            If the target binary has already been instrumented with Edge Coverage tracking instructions, it is executed as is;
            otherwise, in the case of \gls{DBI}, a dynamic translator must dynamically instrument the binary.
        \item \textbf{Coverage Collection}: In the case of Edge Coverage, this part is represented by the data structure used
            to keep track of which edges between \gls{BB}s have been taken and which have not.\\
            This information is used, in some tools such as AFL++ \cite{aflplusplus}, to guide Seed Scheduling,
            a technique that, given different inputs with the same level of coverage (i.e. number of edges visited),
            prioritizes the input that is smaller in size and takes less execution time.
        \item \textbf{Defect Monitor}: This component checks the exit status of the target to catch any crash or
            timeout of the target, or any other unexpected behavior.
            If none of these conditions are met, meaning that the target ends normally, the coverage is checked
            to see whether the current input discovered a new edge.
            If the input discovered a new edge, that input is added to the corpus.
        \end{enumerate}